\section{Trabajos relacionados}

Se comparan al menos tres sistemas o estudios publicados que resuelven un problema equivalente
(gestión de incidencias, mesas de ayuda o sistemas distribuidos de atención). De cada uno se
indica qué se toma y en qué se diferencia el proyecto.

\begin{table}[H]
\centering
\caption{Comparación con trabajos relacionados}
\label{tab:relacionados}
\small
\begin{tabularx}{\linewidth}{@{}clYYY@{}}
\toprule
\# & Obra & Qué resuelve & Qué tomamos & En qué nos diferenciamos \\
\midrule
1 & Li et al. (2022)~\cite{li2022incident}
  & Detección proactiva de incidentes y generación automatizada de tickets ante degradación de
    rendimiento en servicios en la nube a gran escala.
  & Desacoplar el registro inicial del ticket del procesamiento de adjuntos y notificaciones para
    garantizar persistencia inmediata sin bloqueo.
  & Li et al. gestionan incidentes internos de centros de datos y telemetría; SGIRI atiende clientes
    residenciales de ISP, coordinando visitas de técnicos en campo con SLA en horas. \\
\addlinespace
2 & Harsh et al. (2023)~\cite{harsh2023murphy}
  & Diagnóstico y aislamiento de degradaciones intermitentes de rendimiento correlacionando reportes
    en tickets de soporte con trazas distribuidas.
  & Propagación de identificadores de correlación (\texttt{X-Correlation-Id}) en cada petición para
    auditar la trazabilidad completa del ticket entre microservicios.
  & Murphy enfoca el diagnóstico automático de trazas RPC; SGIRI aborda el flujo operativo humano
    completo: asignación por carga, seguimiento móvil y validación de evidencias. \\
\addlinespace
3 & Li et al. (2023)~\cite{li2023identifying}
  & Identificación de cambios de versión o configuración causantes de incidencias reportadas por
    usuarios en sistemas de servicios en línea.
  & Validación estricta de precondiciones y máquina de estados en el límite de la API para rechazar
    transiciones erróneas antes de despachar tareas al personal.
  & Su estudio se centra en incidentes de código en SaaS; SGIRI atiende la logística de técnicos
    desplazados con conectividad móvil intermitente y gestión de cuadrillas locales. \\
\bottomrule
\end{tabularx}
\end{table}
